\chapter{Residues, intersections and cut transfer}
\label{cbc:chap:transfer}

A collision of the first and last positions can cross every cut
while leaving all intermediate positions distinct. Such a coefficient
survives the cut cobar. At three positions it supplies the homotopy
between the two consecutive products. At four positions the crossed
partition $13|24$ supplies a further summand. We construct these
coefficients and their maps before computing their cohomology.

The Heisenberg example then connects two calculations. Wick
contractions compare the coefficients on separated groups, with an
explicit correction to their cut maps. Derived intersections compute
the exterior powers of the primitive module across the diagonals.
The first calculation retains operations. The second retains the
extensions of their coefficient modules.

Throughout, $k$ has characteristic zero and complexes have
cohomological grading. We first use right differential modules on
$X=\mathbb A^1_k$. In the sections on primitive coefficients we
use left modules, as in Chapter~\ref{cbc:chap:heisenberg}. Tensoring
with the canonical bundle converts between the two conventions.
States and normal derivatives are algebraic. All pole orders on
an individual input are finite.

\section{Principal parts with spectators}
\label{cbc:ct:residues}

Let $V$ be a $k$-vector space with translation $T$. Its right
differential module is $\mathscr V=\omega_X\otimes_kV$, with
\[
 (f\,dz\otimes a)\partial_z=-f'\,dz\otimes a-f\,dz\otimes Ta.
\]

For a partition $\pi$ of $[n]$, order the blocks
$B_1,\ldots,B_r$ by their largest labels, called their roots.
Let $i_\pi:X^r\hookrightarrow X^n$ be its diagonal, and let
$j_r$ include the locus of distinct block positions. Put
\begin{equation}\label{cbc:ct:partition}
 \mathcal V_\pi=i_{\pi,*}j_{r,*}j_r^!
                    (\mathscr V^{\boxtimes r}).
\end{equation}
The closed direct image includes every finite normal derivative.
The root order orients its density and the list of block states.
Merging blocks $i<j$ removes root $i$ and retains root $j$.

In a normal coordinate $u$, write $\delta_u=[u^{-1}]$ for the
simple principal part. In the right normal transfer convention,
\begin{equation}\label{cbc:ct:normal}
 \operatorname{pp}_u(u^{-m-1}a)
       =(a\delta_u)\partial_u^m/m!,\qquad m\geq0.
\end{equation}
Regular series have zero principal part. A regular multiplier acts
by its finite Taylor polynomial on any fixed normal derivative.

For example, let $u=z_i-z_j$, let $t=z_e-z_j\ne0$ be a
spectator separation, and let $r,\ell\geq1$. Then
\begin{equation}\label{cbc:ct:spectator}
 (u-t)^{-\ell}[u^{-r}]
  =(-1)^\ell\sum_{a=0}^{r-1}
      \binom{\ell+a-1}{a}t^{-\ell-a}[u^{a-r}].
\end{equation}
This follows by expanding $(1-u/t)^{-\ell}$. Terms with
$a\geq r$ vanish in the principal-part module. Thus taking a
residue retains the spectator poles and only the required normal
jets.

Suppose $V$ is a vertex algebra with vacuum $\mathbf1$, in degree
zero. We use $Y(Ta,u)=\partial_uY(a,u)$ and
$Y(a,u)\mathbf1=e^{uT}a$. For a multiplier $f$ in
\eqref{cbc:ct:partition}, define
\begin{equation}\label{cbc:ct:field-residue}
 r_{ij}(f;a_1,\ldots,a_r)
   =\operatorname{pp}_{z_i-z_j}
       \bigl(f\,Y(a_i,z_i-z_j)a_j\bigr),
\end{equation}
keeping every other state in root order. Extend this formula over
the existing normal transfers. Lower truncation of the field and
the finite pole order of $f$ make its value finite.

The two
translation identities are
\begin{align*}
 r(f;Ta,b)&=\partial_u r(f;a,b)-r(\partial_u f;a,b),\\
 r(f;a,Tb)&=T r(f;a,b)-\partial_u r(f;a,b)
                                      +r(\partial_u f;a,b).
\end{align*}
Here $\partial_u$ differentiates the rational representative.
After tensoring with its density it is minus the right action.
Together with spectator differentiation, they prove that
\eqref{cbc:ct:field-residue} is a map of right differential modules.

\begin{proposition}[The residue relations]\label{cbc:ct:relations}
Disjoint mergers commute. For three blocks $i<j<\ell$,
\begin{equation}\label{cbc:ct:three}
 r_{(ij),\ell}r_{ij}-r_{i,(j\ell)}r_{j\ell}
                         +r_{j,(i\ell)}r_{i\ell}=0.
\end{equation}
The residues commute with restrictions to cut opens and with
surviving exceptional collision restrictions. A vacuum in the
first merged block gives the principal part of the multiplier.
A vacuum in the second gives its principal part multiplied by
the creation series $e^{uT}a$.
\end{proposition}

\begin{proof}
Disjoint mergers use disjoint state factors and independent normal
variables. Their finite principal parts commute. Cuts and surviving
collision restrictions use the same block variables and normal
transfer. The vacuum assertions are the two vacuum field identities.

For three blocks put $s=z_i-z_\ell$ and $t=z_j-z_\ell$.
Mutual locality and lower truncation give a common germ
\[
 m\in V[[s,t]][s^{-1},t^{-1},(s-t)^{-1}]
\]
whose two separated expansions are $Y(a,s)Y(b,t)c$ and
$Y(b,t)Y(a,s)c$. Indeed a sufficiently large power of $s-t$
makes these equal, and their common expansion has bounded negative
powers of both $s$ and $t$. Weak associativity identifies its
expansion near $s=t$ with $Y(Y(a,s-t)b,t)c$.

Write $P_{a,b}$ for the oriented iterated principal part along
$a=0$, then $b=0$. The three-channel identity for such a germ is
\[
 P_{s-t,t}-P_{t,s}+P_{s,t}=0.
\]
It is the boundary squared in the localization complex of the
three lines. Partial fractions in $s$ over $k((t))$ give a
second verification. The poles at $0$ and $t$ give the two
terms in the expansion at infinity. The polynomial part has
zero second principal part.

Multiplication by powers of $s,t,s-t$ and differentiation
preserve the identity, so it holds with every finite normal
derivative. Transport the density once in each change of normal
coordinates. Applying this identity to $fm$ proves
\eqref{cbc:ct:three} with its displayed signs.
\end{proof}

For the construction below, a \emph{residue system} means maps
on \eqref{cbc:ct:partition} satisfying
Proposition~\ref{cbc:ct:relations}. This permits the same
construction whenever these coefficient identities have been
proved, independently of a presentation by fields.

\section{The finite first-cut contraction}
\label{cbc:ct:contraction}

A cut $c\in\{1,\ldots,n-1\}$ is \emph{admissible} for
$\pi$ if each block lies on one side of $c$. Write $J(\pi)$
for the set of admissible cuts. On the cut open $U_c$ of
\eqref{eq:words-cut-product}, a partition crossing $c$ has zero
restriction. An admissible partition splits into its left and
right coefficients. Since its distinct-block locus already lies
in $U_c$, its coefficient is the direct image from that open.

Give the partition sum the shifts and differential
\begin{equation}\label{cbc:ct:coalgebra}
 C_n=\bigoplus_\pi\mathcal V_\pi[|\pi|],\qquad
 d_C|_\pi=\sum_{i<j}(-1)^{i-1}r_{ij},\qquad C_0=k.
\end{equation}
The two orders of a disjoint merger have opposite signs. The
three orders reaching one triple block give
\eqref{cbc:ct:three}. Thus $d_C^2=0$.

Restriction and splitting
on every $U_c$ give a strong cut coproduct. The differential on
the two sides is the tensor differential. Successive cuts
restrict to the same open set.

The positive cut cobar has one cell $(\pi,Q)$ for each
$Q\subseteq J(\pi)$:
\begin{equation}\label{cbc:ct:cells}
 K_n=\bigoplus_{\pi,Q\subseteq J(\pi)}
       \mathcal V_\pi[|\pi|-|Q|-1].
\end{equation}
For states in degree zero its cell degree is $|Q|+1-|\pi|$.

Write $D=v+\delta$, where $v$ inserts a cut and $\delta$
merges two blocks in one of the segments of $Q$. If $r_{\leq c}$
counts blocks before $c$, and $i<j$ are global block indices in
segment $\ell$, the signs are
\begin{align}
 v_c&=(-1)^{\#\{q\in Q:q<c\}+1-r_{\leq c}}\operatorname{id},
                    &&c\in J(\pi)\setminus Q,
                    \label{cbc:ct:cut-sign}\\
 \delta_{ij}&=(-1)^{\ell-1+i}r_{ij}.
                    \label{cbc:ct:merge-sign}
\end{align}

These follow from the cobar rule
$d(tc)=-t(d_Cc)+\sum(-1)^{|c'|+1}(tc')(tc'')$,
where $t$ has degree one and $r_{\mathrm{before}}$ counts
the blocks in the segments preceding $\ell$.
Moving past the preceding segments
contributes $\ell-1-r_{\mathrm{before}}$. The internal merger
contributes $i-r_{\mathrm{before}}$. The split contributes
$1-r_{\leq c}+r_{\mathrm{before}}$. This gives both signs.
The residue relations and cut compatibility give
$v^2=\delta^2=v\delta+\delta v=0$.

A partition is \emph{indecomposable} when $J(\pi)$ is empty.
A merger cannot create an admissible cut, so the surviving cells
form the subcomplex
\begin{equation}\label{cbc:ct:survivors}
 G_n=\bigoplus_{J(\pi)=\varnothing}\mathcal V_\pi[|\pi|-1],
       \qquad d_G|_\pi=\sum_{i<j}(-1)^ir_{ij}.
\end{equation}
Let $i:G_n\to K_n$ be the inclusion and $p_0$ the summand
projection. If $c_0=\min J(\pi)$ exists, set
\begin{equation}\label{cbc:ct:first-cut}
 h_0(\pi,Q)=
 \begin{cases}
  (-1)^{1-r_{\leq c_0}}(\pi,Q\setminus\{c_0\}),&c_0\in Q,\\
  0,&c_0\notin Q.
 \end{cases}
\end{equation}
On indecomposable partitions set $h_0=0$.

\begin{theorem}[Finite residue retract]\label{cbc:ct:retract}
The finite formulas
\begin{equation}\label{cbc:ct:finite-retract}
 \Theta=\sum_{a=0}^{n-1}(-\delta h_0)^a,
       \qquad \rho=p_0\Theta,\qquad H=h_0\Theta
\end{equation}
give maps of differential-module complexes and a homotopy
\[
 \rho D=d_G\rho,\quad D i=i\,d_G,\quad \rho i=1,\quad
 DH+HD=1-i\rho,\quad H^2=Hi=\rho H=0.
\]
They commute with monotone collision restrictions under the
coefficient identifications of Chapter~\ref{chap:collision-words}.
\end{theorem}

\begin{proof}
Adding and removing $c_0$ in either order gives the identity.
The terms at any other cut cancel, since its position in the
ordered cut list changes by one when $c_0$ is removed. Hence
\[
 vh_0+h_0v=1-ip_0,\qquad
 h_0^2=h_0i=p_0h_0=0,\qquad p_0i=1.
\]
Each $\delta$ lowers the number of blocks, so
$(\delta h_0)^n=0$ and $\Theta=(1+\delta h_0)^{-1}$.

Put $E=1+\delta h_0$. Using $\delta i=i\,d_G$ gives
\[
 DE-Ev=\delta-\delta(vh_0+h_0v)=\delta ip_0.
\]
Since $h_0\delta i=0$, one has $\Theta\delta i=\delta i$.
Multiplication on both sides by $\Theta$ therefore gives
$\Theta D=v\Theta+i\,d_Gp_0\Theta$. Projecting proves the
chain-map identity for $\rho$. It also gives
\[
 DH+HD=(vh_0+h_0v+\delta h_0)\Theta=1-i\rho.
\]

Finally $\Theta i=i$ and $\Theta h_0=h_0$ give all side
conditions. Every term is a finite composition of coefficient
maps and keeps its full normal transfer.

Under a monotone collision only partitions whose blocks are
unions of collision fibers survive. No admissible cut of such a
partition lies inside a fiber. Its cuts therefore identify with
those of the quotient partition, with the same order and block
counts. This preserves $h_0$, its signs and the residues. The
identifications compose because they preserve the same lists.
\end{proof}

At two positions,
\[
 K_2=[\mathcal V_{1|2}\xrightarrow{(-r_{12},1)}
           \mathcal V_{12}\oplus\mathcal V_{1|2}],\qquad
 \rho(p,q)=p+r_{12}q,
\]
in degrees $-1,0$. The residue term in $\rho$ is necessary.

At three positions,
\[
 G_3=[\mathcal V_{2|13}\xrightarrow{-r_{2,(13)}}
                         \mathcal V_{123}].
\]
The first term has degree $-1$.

At four positions the terms,
in increasing degree, are
\begin{equation}\label{cbc:ct:four}
 \mathcal V_{2|3|14}
 \xrightarrow{\left(\begin{smallmatrix}
 -r_{2,14}\\r_{3,14}\\-r_{2,3}\\0
 \end{smallmatrix}\right)}
 \begin{matrix}
 \mathcal V_{3|124}\oplus\mathcal V_{2|134}\\
 {}\oplus\mathcal V_{23|14}\oplus\mathcal V_{13|24}
 \end{matrix}
 \xrightarrow{\ (-r_{3,124},-r_{2,134},-r_{23,14},-r_{13,24})\ }
 \mathcal V_{1234}.
\end{equation}
Their degrees are $-2,-1,0$. The composite is
$r_{3,124}r_{2,14}-r_{2,134}r_{3,14}+r_{23,14}r_{2,3}=0$.
The crossed summand $13|24$ has its own outgoing residue even
though the degree-$-2$ term does not map to it.

\section{Transferred products}
\label{cbc:ct:products}

Put $K=(K_n)_{n\geq1}$ and $G=(G_n)_{n\geq1}$.
Concatenation of cobar words inserts the cut between them. It
gives the associative product $\mu$ on $K$, with the empty unit
adjoined in arity zero. Let $s:K\to K[1]$ have degree $-1$.
Set
\[
 b_1(sx)=-sDx,\qquad b_2(sx,sy)=(-1)^{|x|}s\mu(x,y),
 \qquad B_1(sx)=-sd_Gx.
\]

Write $\bar i=sis^{-1}$, $\bar\rho=s\rho s^{-1}$ and
$\bar H=-sHs^{-1}$. Define
\begin{equation}\label{cbc:ct:trees}
 F_1=\bar i,\qquad
 L_r=\sum_{a+b=r}b_2(F_a\otimes F_b),\qquad
 F_r=-\bar H L_r,\qquad B_r=\bar\rho L_r\quad(r\geq2).
\end{equation}
Tensor composition includes the Koszul signs. Each $F_r$ has
degree zero and each $B_r$ degree one.

\begin{theorem}\label{cbc:ct:transfer-products}
The maps $B_r$ define an $A_\infty$-algebra on $G$ for the
cut product. The maps $F_r$ define an $A_\infty$
quasi-isomorphism $G\to K$. All sums are finite in each arity.
On singleton inputs,
\begin{equation}\label{cbc:ct:outer}
 B_n(sx_1,\ldots,sx_n)=s\,r_{1n}(x_1,\ldots,x_n),
                \qquad n\geq2.
\end{equation}
Its target is the coefficient of
$2|\cdots|n-1|\{1,n\}$ in degree $2-n$ of $G_n$.
\end{theorem}

\begin{proof}
Let $T^c_{\mathrm{cut}}$ be the reduced tensor coalgebra for
the cut product, with deconcatenation. Extend $b_1$ to
$T^c_{\mathrm{cut}}(sK)$, writing the resulting differential
as $d$, and extend $b_2$ as a coderivation
$\beta$. The dg associative identities say
$d\beta+\beta d=\beta^2=0$. The tensor extension of the
retract has maps $I_0=T^c_{\mathrm{cut}}(\bar i)$,
$P_0=T^c_{\mathrm{cut}}(\bar\rho)$ and
homotopy, on length $a$,
\[
 H_0=\sum_{j=0}^{a-1}(\bar i\bar\rho)^{\otimes j}
           \otimes\bar H\otimes1^{\otimes(a-j-1)}.
\]
The terms telescope to $dH_0+H_0d=1-I_0P_0$, with the
three zero side conditions.

Define
\[
 Q=\beta(1+H_0\beta)^{-1},\quad
 I=(1+H_0\beta)^{-1}I_0,\quad
 P=P_0(1+\beta H_0)^{-1}.
\]
Every inverse is a finite geometric series, since $\beta$
lowers tensor length. Substituting
$Q=\beta-\beta H_0Q$ into the contraction identity gives
\[
 dQ+Qd=-QI_0P_0Q.
\]
It follows by multiplication that
$B=B_1+P_0QI_0$ squares to zero,
$(d+\beta)I=IB$, $P(d+\beta)=BP$ and $PI=1$.

For deconcatenation $\Delta$ one has
$\Delta H_0=(H_0\otimes1+I_0P_0\otimes H_0)\Delta$.
Apply this identity to $I=I_0-H_0\beta I$ and
$P=P_0-P\beta H_0$. Induction on tensor length, using
$H_0I=PH_0=0$, shows that $I,P$ are coalgebra maps.
Consequently $B=P(d+\beta)I$ is a coderivation. Its
length-one components are exactly \eqref{cbc:ct:trees}.
These facts prove the $A_\infty$ identities and the
quasi-isomorphism assertion.

For singleton inputs let $e_n$ denote the uncut cell and $w_c$
the cell with sole cut $c$. The product and suspension signs
give, inductively,
\[
 L_n=\sum_{c=1}^{n-1}(-1)^{1-c}s w_c,
       \qquad Hw_1=e_n,\qquad Hw_c=0\ (c>1).
\]
After removing the first cut, there is no cut left for another
$h_0$. Thus $F_n=se_n$.

Also $\rho w_c=0$ for $c>1$,
whereas $\rho w_1=-p_0\delta e_n$. Of the single mergers
of the singleton partition, only $1\sim n$ leaves no admissible
cut. Its sign in $\delta$ is negative. This proves
\eqref{cbc:ct:outer}.
\end{proof}

\begin{corollary}[The coefficients of higher mergers]
\label{cbc:ct:merger-signs}
Let $m\geq2$, and take inputs $x_a$ with indecomposable partitions
$\pi_1,\ldots,\pi_m$ on consecutive groups of positions.
Put $r_a=|\pi_a|$ and $r=\sum_a r_a$. Number the blocks of
their union in root order. For state coefficients in degree zero,
\begin{equation}\label{cbc:ct:suspended-mergers}
 B_m(sx_1,\ldots,sx_m)
 =s\sum_{i=1}^{r_1}\sum_{j=r-r_m+1}^{r}
                    (-1)^{i+1}r_{ij}(x_1,\ldots,x_m).
\end{equation}
Here $r_{ij}$ acts on the union coefficient.

Define the desuspended map $m_m=s^{-1}B_m s^{\otimes m}$.
Since $|x_a|=1-r_a$, its formula is
\begin{equation}\label{cbc:ct:desuspended-mergers}
 m_m(x_1,\ldots,x_m)
 =\sum_{i=1}^{r_1}\sum_{j=r-r_m+1}^{r}
 (-1)^{i+1+\sum_{a=1}^m(m-a)(1-r_a)}
                         r_{ij}(x_1,\ldots,x_m).
\end{equation}
Every formal merger coefficient for this contraction is $1$ or $-1$.
\end{corollary}

\begin{proof}
The union partition has admissible cuts exactly at the input
boundaries. We prove by induction that $F_a$ sends its suspended
inputs to the suspended uncut union with coefficient one.
This holds for $F_1$.

A split $p+q=a$ in \eqref{cbc:ct:trees} creates a cell with
one cut, after input $p$. On such a cell $H=h_0$.
Indeed, removing its cut leaves no cut for the next $h_0$,
and a residue does not insert one. Only $p=1$ survives, because
$h_0$ removes the first admissible cut.
The signs of $b_2$ and $h_0$ are both $(-1)^{1-r_1}$.
They cancel, proving the induction.

At the root, the only part of $\rho$ that survives on a
one-cut cell is $-p_0\delta h_0$. Again only $p=1$ contributes.
The product and homotopy signs cancel as above.
The merger $r_{ij}$ then has coefficient $(-1)^{i+1}$.
A single merger removes every input boundary exactly when it
joins the first input to the last. Each such merger survives
$p_0$ once, proving \eqref{cbc:ct:suspended-mergers}.

Evaluating $s^{\otimes m}$ on homogeneous inputs contributes
the sign $(-1)^{\sum_a(m-a)|x_a|}$ by the tensor Koszul rule.
This proves \eqref{cbc:ct:desuspended-mergers}.
\end{proof}

For example, on inputs with partitions $2|13$, $4$ and $5$,
the desuspended ternary map is $r_{2,5}-r_{13,5}$.
The formula fixes the first-cut contraction. It makes no claim
about coefficients obtained from a different contraction.

On three singleton inputs, Theorem~\ref{cbc:ct:transfer-products}
gives the homotopy $h=r_{13}$ and
\[
 d_Gh=r_{(12),3}r_{12}-r_{1,(23)}r_{23}.
\]
In arity four the expression
$L_4=b_2(F_1,F_3)+b_2(F_2,F_2)+b_2(F_3,F_1)$ contains
the five planar binary trees. The square-zero identity for $B$
includes all their boundary terms with the Koszul signs. Its
maps are defined on the common localized coefficient. They use
no identification of incompatible iterated Laurent series.

The construction also applies to dg state coefficients when
residues are cochain maps and all permutations use the Koszul
rule. On a cell shifted by $r-|Q|-1$, the internal differential
is multiplied by $(-1)^{r-|Q|-1}$. Tensoring the first-cut
contraction with the state complexes gives the same proof.

For a concrete three-position value, take the Heisenberg current
with $[J_m,J_n]=\kappa m\delta_{m+n,0}$, so
$J=J_{-1}\mathbf1$ and $TJ=J_{-2}\mathbf1$. Set
$s=z_1-z_3$, $t=z_2-z_3$ and use the multiplier $t^{-2}$
on three current inputs. Let $L,R,N$ first merge $12,23,13$,
respectively, and then merge the remaining pair. In the full
normal module at $s=t=0$ their values are
\begin{align*}
 L&=\kappa(s^{-2}t^{-2}+2s^{-3}t^{-1})J,\\
 R&=\kappa(2s^{-2}t^{-2}+2s^{-3}t^{-1})J
                         +\kappa s^{-2}t^{-1}TJ,\\
 N&=\kappa s^{-2}t^{-2}J+\kappa s^{-2}t^{-1}TJ.
\end{align*}

These are principal-part classes, with their density understood.
Indeed
\[
 \operatorname{pp}_t(t^{-2}Y(J,t)J)
    =\kappa t^{-4}\mathbf1+t^{-2}J^2+t^{-1}(TJ)J.
\]

Contraction of the remaining current with $J^2$ gives
$2\kappa s^{-2}J$. Contraction with $(TJ)J$ gives
$2\kappa s^{-3}J+\kappa s^{-2}TJ$. This proves the value
of $R$. The creation series gives $N$.

For $L$, use the
normal transfer
\[
 [(s-t)^{-2}t^{-2}]=[s^{-2}t^{-2}]+2[s^{-3}t^{-1}],
\]
obtained by differentiating the simple double principal part.
Thus $L-R+N=0$ and $N\ne0$. Removing the nonconsecutive
channel would remove the required homotopy while leaving a
nonzero associativity defect.

\section{Wick comparisons across a cut}
\label{cbc:ct:wick}

Two different comparisons arise for polynomial Heisenberg
coefficients. The first changes the field kernel. The second
contracts a Wick operator by a homotopy of the primitive
complex. We specify their domains separately.

Use polynomial states $V=k[x_1,x_2,\ldots]$ with
$Tx_p=px_{p+1}$, vacuum $1$ and current $J=x_1$. For an even
Laurent polynomial $P(u)$, put
\begin{equation}\label{cbc:ct:kernel}
 K^P_{pq}(u)=\frac{(-1)^{q-1}P^{(p+q-2)}(u)}
                         {(p-1)!(q-1)!},\qquad
 C_P(u)=\sum_{p,q\geq1}K^P_{pq}(u)
                         \partial_{x_p}\otimes\partial_{x_q}.
\end{equation}
With $m$ the ordinary polynomial multiplication, the fields are
$Y_P(u)=m(e^{uT}\otimes1)e^{C_P(u)}$.

Only finitely many derivatives act on each polynomial, and each
contraction lowers polynomial degree by two. The exponential
therefore terminates. The recurrences
$pK^P_{p+1,q}=(K^P_{pq})'$ and
$qK^P_{p,q+1}=-(K^P_{pq})'$ give translation covariance.
Commuting contractions give the common three-field expansion
with exponent $C_{P,12}+C_{P,13}+C_{P,23}$. Evenness gives
mutual locality after clearing its finitely many poles.
Thus these fields give the residue system just constructed.

For two such kernels put $H=P-Q$. On a partition coefficient
with distinct block roots define
\[
 \Gamma^H_\pi=\exp\left(\sum_{a<b}C_H(z_a-z_b)_{ab}\right).
\]
Extend this horizontal operator to the closed direct image
in \eqref{cbc:ct:partition}.
For two separated groups $A,B$ in their word order, let
$C_{H;A,B}$ be the sum of contractions between the groups.
Write $\Delta^P_{A,B}$ for the cut restriction and splitting
in the coefficient system of $Y_P$.

\begin{proposition}\label{cbc:ct:wick-cuts}
The operators $\Gamma^H$ are horizontal isomorphisms of the
collision complexes, preserve vacuum coefficients, and satisfy
\begin{align}
 \Gamma^H_{\pi/(i\sim j)}r^P_{ij}
       &=r^Q_{ij}\Gamma^H_\pi,\label{cbc:ct:wick-contact}\\
 (\Gamma^H_A\boxtimes\Gamma^H_B)\Delta^P_{A,B}
                    (\Gamma^H_{A\cup B})^{-1}
       &=e^{-C_{H;A,B}}\Delta^Q_{A,B}.
                                  \label{cbc:ct:wick-cut}
\end{align}
The cross factors obey the three-group cocycle identity and
commute with mergers inside either group.
\end{proposition}

\begin{proof}
The recurrences in \eqref{cbc:ct:kernel} give
$[\partial_{z_a}+T_a,C_{H,ab}]=0$ and the corresponding
identity at $b$, so $\Gamma^H$ is horizontal. At a merger,
a derivative of the merged polynomial splits into its derivatives
on the two inputs. Taylor translation changes each spectator
separation by the collision coordinate. The internal contraction
satisfies $Y_P=Y_Qe^{C_H}$. These identities, followed by the
finite Taylor rule \eqref{cbc:ct:spectator}, prove
\eqref{cbc:ct:wick-contact}.

Evenness makes the same argument
valid after permuting blocks into adjacent positions. It also
preserves the root-change transfers. Vacuum derivatives vanish.

Split the exponent of $\Gamma^H_{A\cup B}$ into its two
internal sums and $C_{H;A,B}$. All contractions commute.
Cancelling the internal sums gives \eqref{cbc:ct:wick-cut}.
For three groups both iterated cross factors have exponent
$-C_{H;A,B}-C_{H;A,C}-C_{H;B,C}$. The merger rule already
proved gives compatibility within each group.
\end{proof}

For $Q=0$, applying $(\Gamma^P)^{-1}$ to regular polynomial
tensors gives residue-free sections, including
\[
 J_1J_2-P_{12},\qquad
 J_1J_2J_3-P_{12}J_3-P_{13}J_2-P_{23}J_1,
 \qquad P_{ij}=P(z_i-z_j).
\]
The zero-kernel residue-free lattice itself contains polar
sections: $u^{-1}(J\otimes1-1\otimes J)$ fuses to
$(e^{uT}J-J)/u$, which is regular. It is therefore larger than
the unlocalized external product. Also $\Gamma_1^H=1$
while $\Gamma_2^H(J\otimes J)=J\otimes J+H(u)1\otimes1$.
The coefficient comparison is generally not the tensor power
of a state map.

The primitive complex permits a homotopy of the cross
contraction itself. Fix $\kappa\in k^\times$ and use
$K_{\kappa,n}$ of \eqref{cbc:hc:linear-complex} and its
polynomial algebra $Z_{\kappa,n}$, with
$dx(f)=-\kappa\gamma(f')$ and $d\gamma(f)=0$.
On an open separating groups $s,t$, the principal parts split
by groups by Proposition~\ref{cbc:hc:cuts}.

For proper rational
$f$ in group $s$ and $g$ in group $t$, let
$\operatorname{Res}_{D_t}$ sum residues at the whole divisor
of group $t$. Define a degree-zero contraction $\Delta_{st}$
by its pair value
\begin{equation}\label{cbc:ct:wick-pair}
 \Delta_{st}(x(f)x(g))
       =\kappa\operatorname{Res}_{D_t}g\,df,
\end{equation}
and zero on pairs involving a ghost. It acts by summing over
unordered pairs of factors in a polynomial, with the Koszul
sign. Define degree-$-1$ pair contractions by
\begin{equation}\label{cbc:ct:wick-primitives}
 K^L_{st}(\gamma(f)x(g))=-\operatorname{Res}_{D_t}fg\,dw,
 \qquad
 K^R_{st}(x(f)\gamma(g))=-\operatorname{Res}_{D_s}fg\,dw,
\end{equation}
and zero on the other oriented ghost pair.

\begin{theorem}[Polynomial homotopies for Wick contraction]
\label{cbc:ct:wick-homotopy}
On this separated-group open,
$[d,K^L_{st}]=[d,K^R_{st}]=\Delta_{st}$.
The contractions commute with the position connections and
internal collisions. For finitely many group pairs $e$, choose
$K_e$ to be either primitive or their average. Put
$\Omega=k[t_e,dt_e]_e$, with $|t_e|=0$, $|dt_e|=1$ and
$d_\Omega t_e=dt_e$. With these polynomial forms written
on the left, the map
\begin{equation}\label{cbc:ct:wick-cube}
 \mathcal F=\exp\left(\sum_e t_e\Delta_e+dt_eK_e\right)
\end{equation}
from $Z_{\kappa,n}$ to $\Omega\otimes_kZ_{\kappa,n}$ is a
cochain map. Its faces and their integrals give compatible
Wick comparisons and homotopies for refinements of the groups.
\end{theorem}

\begin{proof}
On $x(f)x(g)$, the left anticommutator is
$\kappa\operatorname{Res}_{D_t}gf'\,dw$. The right one is
$\kappa\operatorname{Res}_{D_s}fg'\,dw$. Integration by
parts identifies them: the proper rational differential $fg'\,dw$
has no residue at infinity, so the sum over the two groups is
zero. On pairs with a ghost both anticommutators vanish.
Pair values determine these constant contraction operators.

Group residue is computed by monic division. It commutes with
parameter differentiation and with internal collision substitution,
without inverting differences inside the group. Thus the operators
preserve these maps. Constant graded derivatives supercommute.

Since forms stand on the left,
\[
 [d+d_{\mathbf t},t_e\Delta_e+dt_eK_e]
     =dt_e\Delta_e-dt_e[d,K_e]=0.
\]
The degree-zero exponent terms commute and lower polynomial
degree. Their finite exponential is a cochain map. Restriction
to faces gives the endpoint maps. Polynomial Stokes gives their
homotopy identities. Refinement partitions the same set of cross
pairs, so all its comparisons are faces of this formula.
\end{proof}

For $\Delta=\sum_e\Delta_e$ and $K=\sum_eK_e$, the
one-parameter integral is
\[
 S=K\sum_{j\geq0}\frac{\Delta^j}{(j+1)!},\qquad
       [d,S]=e^\Delta-1.
\]
This sum terminates on every polynomial. For
$f=(w-z_1)^{-m}$, $g=(w-z_2)^{-n}$ and $u=z_1-z_2$,
the pair values are
\begin{align}
 \beta_{mn}&=\kappa(-1)^m
       \frac{(m+n-1)!}{(m-1)!(n-1)!}u^{-m-n},
                                     \label{cbc:ct:wick-values}\\
 K^L_{mn}&=(-1)^{m+1}\binom{m+n-2}{n-1}u^{1-m-n},\qquad
 K^R_{mn}=(-1)^m\binom{m+n-2}{m-1}u^{1-m-n}.
                                     \notag
\end{align}

These follow by Taylor expansion at the indicated group's pole.
At $m=n=1$ the primitives are $u^{-1}$ and $-u^{-1}$,
and the contraction is $-\kappa u^{-2}$. The sign uses the
rational-generator convention $dx(f)=-\kappa\gamma(f')$.
Equation~\eqref{cbc:ct:kernel} uses the current convention
with $Y_P(J,u)J=P(u)+(e^{uT}J)J$.

The polynomial cohomology map of
Theorem~\ref{cbc:hc:polynomial-cohomology} can be composed with
these homotopies. Their domain allows every internal collision
but keeps different groups disjoint. Extension through an external
collision with both vacuum insertions remains an additional
comparison problem.

\section{Derived intersections of diagonal coefficients}
\label{cbc:ct:intersections}

The diagonal filtration of $H_n$ in
Theorem~\ref{cbc:hc:filtration} has nonflat quotients. Already
for $R=k[u]$ and $P=R[u^{-1}]/R$, every tensor $p\otimes q$
vanishes: some $u^m$ kills $p$, while $q=u^mq'$ for a
principal part $q'$. Hence $P\otimes_RP=0$. The flat
resolution $[R\to R[u^{-1}]]$ in degrees $-1,0$ instead
gives
\[
 P\otimes_R^{\mathrm L}P\simeq P[1].
\]
This degree is needed to pass from the diagonal filtration to
exterior powers.

Continue over $R_n=k[z_1,\ldots,z_n]$ with its left Weyl
algebra $D_n$. For a partition $\pi$, let $\Delta_\pi$ be
its diagonal and put $c_\pi=n-|\pi|$. Write
\[
 A_\pi=R\Gamma_{\Delta_\pi}R_n,\qquad
 \delta_\pi=H^{c_\pi}_{\Delta_\pi}(R_n).
\]
Choose an ordered list of independent linear normal equations
$f_1,\ldots,f_{c_\pi}$. Then
\begin{equation}\label{cbc:ct:localization}
 A_\pi\simeq\bigotimes_{a=1}^{c_\pi}
     [R_n\longrightarrow R_n[f_a^{-1}]]
       \simeq\delta_\pi[-c_\pi],
\end{equation}
where each bracket has degrees zero and one.

The localizations
are flat and carry their derivative connections. The cohomology
is concentrated in degree $c_\pi$: after a linear coordinate
change this is the tensor product of the one-variable
principal-part calculation. Its Weyl module is the trivial
module on the diagonal, pushed forward with all normal
derivatives.

\begin{proposition}\label{cbc:ct:diagonal-tensor}
Let $\pi\vee\sigma$ impose both sets of equalities. There is
an equivalence of differential-module complexes
\begin{equation}\label{cbc:ct:excess}
 \delta_\pi\otimes_{R_n}^{\mathrm L}\delta_\sigma
 \simeq\delta_{\pi\vee\sigma}
            [c_\pi+c_\sigma-c_{\pi\vee\sigma}].
\end{equation}
The localization complexes retain the symmetric tensor
coherence. Ordered normal coordinates specify the signs in
their top-principal-part identifications.
\end{proposition}

\begin{proof}
Tensoring the localization complexes concatenates the lists of
normal equations. This computes local cohomology on the
intersection of their supports, so
$A_\pi\otimes^{\mathrm L}A_\sigma\simeq A_{\pi\vee\sigma}$.
Restoring the shifts in \eqref{cbc:ct:localization} gives
\eqref{cbc:ct:excess}. Localization, its counit and their tensor
maps commute with differentiation. Permutations of equations
act by the Koszul signs. A change of independent normal basis
uses the corresponding determinant in the principal-part
identification.
\end{proof}

For a graph $G$ on $[n]$, allow repeated edges but no loops,
and let $\pi_G$ be its connected-component partition. Repeated
application of \eqref{cbc:ct:excess} gives
\begin{equation}\label{cbc:ct:graph}
 \bigotimes_{e\in E(G)}^{\mathrm L}\delta_e
  \simeq\delta_{\pi_G}[\beta_1(G)],\qquad
 \beta_1(G)=|E(G)|-n+|\pi_G|.
\end{equation}

A spanning forest supplies independent normals. Each remaining
edge supplies one shift. For example, with $a=z_1-z_2$ and
$b=z_2-z_3$ and $p,q\geq1$, every $[a^{-p}b^{-q}]$ is killed by
$(a+b)^{p+q-1}$. Localizing at the third edge therefore gives
zero, and the triangle tensor is $\delta_{123}[1]$.
The coherent structure sheaf $R_n/I_\pi$, where $I_\pi$
defines the diagonal, is a different module
and has different derived intersections.

In characteristic zero, the derived exterior power
$\mathrm L\Lambda^rM$ is the image of
$r!^{-1}\sum_{\sigma\in S_r}\operatorname{sgn}(\sigma)\sigma$
on $M^{\otimes^{\mathrm L}r}$. Permutations here include
the Koszul signs of the complex.

\begin{lemma}\label{cbc:ct:normal-exterior}
If $\delta=H^c_Z(R_n)$ for a linear diagonal $Z$ and $r>1$,
then
\[
 \mathrm L\Lambda^r\delta\simeq
 \begin{cases}
  0,&c\text{ even},\\
  \delta[(r-1)c],&c\text{ odd}.
 \end{cases}
\]
For $r=1$ the exterior power is $\delta$.
\end{lemma}

\begin{proof}
Put $A=R\Gamma_ZR_n$. Its counit $A\to R_n$ gives a
symmetric equivalence $A^{\otimes r}\simeq A$. Indeed the
tensor is supported on $Z$. Its map to $R_n$ lifts uniquely,
in the derived mapping space, through the right adjoint
$R\Gamma_Z$. Applying the counit to all but one factor gives
equivalences with this same lift. This supplies their symmetric
coherence.

Since $\delta=A[c]$, a permutation of the shifted factors
introduces the character $\operatorname{sgn}^c$. The alternating
idempotent vanishes for even $c$ and is the identity for odd
$c$. In the latter case the shift is $rc-c$. In particular
the flip on $P\otimes^{\mathrm L}P\simeq P[1]$ is minus
the identity.
\end{proof}

\section{Exterior powers of the primitive coefficient}
\label{cbc:ct:exterior}

Let $S(b,j)$ be the number of partitions of a $b$-element set
into $j$ nonempty blocks. For $b\geq2$ define
\[
 Q_b(t)=\sum_{j=1}^b S(b,j)(j-1)!t^j,
 \qquad s_1(\pi)=\#\{B\in\pi:|B|=1\}.
\]
For example, $Q_2=t+t^2$ and $Q_3=t+3t^2+2t^3$.

\begin{theorem}[Diagonal multiplicities]\label{cbc:ct:multiplicities}
For $q\geq0$, the exterior power $\Lambda^q_{R_n}H_n$ is flat over $R_n$
and has finite length over $D_n$. It vanishes for $q>n$.
The composition multiplicity of $\delta_\pi$ is
\begin{equation}\label{cbc:ct:multiplicity}
 m_{\pi,q}=[t^q](1+t)^{s_1(\pi)}
                 \prod_{\substack{B\in\pi\\|B|\geq2}}
                      Q_{|B|}(t).
\end{equation}
Consequently
\begin{equation}\label{cbc:ct:length}
 \operatorname{length}_{D_n}\Lambda^qH_n=q!S(n+1,q+1).
\end{equation}
These are composition multiplicities. The resulting modules
need not split into their simple constituents.
\end{theorem}

\begin{proof}
An exterior power is a direct summand of the corresponding
tensor power, hence is flat by Corollary~\ref{cbc:hc:flat}.
It is torsion-free over the domain $R_n$. Its generic fiber
is the exterior power of the $n$-dimensional residue space,
which proves vanishing above $n$.

Consider the Grothendieck group of bounded differential-module
complexes whose cohomology has finite filtrations by the
$\delta_\pi$. These diagonal modules are simple. On a finite
sum, polynomials in the coordinate Euler operators isolate a
nonzero monomial. Tangential differentiation and normal
multiplication then reduce it to the simple normal principal
part. The reverse operations generate every monomial.

Modules with
different supports are nonisomorphic. Thus their classes are
independent. Derived tensor remains in this category by
Proposition~\ref{cbc:ct:diagonal-tensor}.

Set $e_\pi=(-1)^{c_\pi}[\delta_\pi]$. Their multiplication
is $e_\pi e_\sigma=e_{\pi\vee\sigma}$. Write
$\lambda_t(M)=\sum_{r\geq0}[\mathrm L\Lambda^rM]t^r$.
For an exact triangle $A\to B\to C$, the finite tensor
filtration, followed by the alternating idempotent, gives
$\lambda_t(B)=\lambda_t(A)\lambda_t(C)$. The preceding lemma
therefore gives
\[
 \lambda_t(\delta_\pi)=
 \begin{cases}
  1+e_\pi t,&c_\pi\text{ even},\\
  (1+e_\pi t)^{-1},&c_\pi\text{ odd}.
 \end{cases}
\]
Each coefficient of this formal series involves only a finite
derived calculation.

For a partition $\tau$, evaluate $e_\pi$ as one when $\pi$
refines $\tau$ and as zero otherwise. These are ring
homomorphisms, and the refinement matrix is triangular with
diagonal entries one. The subset filtration of $H_n$ gives
\[
 \operatorname{ev}_\tau\lambda_t(H_n)
 =(1+t)^{\sum_{B\in\tau}\sum_{s=1}^{|B|}
                        (-1)^{s-1}\binom{|B|}{s}}
 =(1+t)^{|\tau|}.
\]
The $n$ singleton constituents are separate copies of the same
regular module. The exponent counts all of them. These
evaluations are operations on classes, not fibers of modules.

Write $\lambda_t(H_n)=\sum_\pi a_\pi(t)e_\pi$.
M\"obius inversion in the refinement order gives
\[
 a_\pi(t)=\prod_{B\in\pi}p_{|B|}(1+t),\qquad
 p_b(x)=\sum_{j=1}^bS(b,j)(-1)^{j-1}(j-1)!x^j.
\]

To see the coefficient, a partition of one block into $j$
pieces has M\"obius value $(-1)^{j-1}(j-1)!$ when those
pieces are joined. Values multiply on different blocks. Their
generating series is
\[
 \sum_{b\geq1}p_b(x)\frac{z^b}{b!}
       =\log\bigl(1+x(e^z-1)\bigr).
\]
At $x=1+t$ this becomes
$z+\log(1+t(1-e^{-z}))$. Hence $p_1(1+t)=1+t$ and
$(-1)^{b-1}p_b(1+t)=Q_b(t)$ for $b\geq2$.
The sign changing $e_\pi$ back to $[\delta_\pi]$ cancels
these block signs and proves \eqref{cbc:ct:multiplicity}.

The finite derived tensor filtrations have finitely many
diagonal simple constituents in each exterior degree. Flatness
identifies the derived exterior power with the ordinary one in
degree zero. The computed coefficients are therefore actual
composition multiplicities. Summing over partitions, with the
empty-set term included, gives
\[
 \sum_{n\geq0}\sum_q
  \operatorname{length}\Lambda^qH_n\,t^q\frac{z^n}{n!}
     =\frac{e^z}{1-t(e^z-1)}.
\]
The coefficient of $t^q$ is
$e^z(e^z-1)^q=(q+1)^{-1}\partial_z(e^z-1)^{q+1}$.
Its coefficient of $z^n/n!$ is $q!S(n+1,q+1)$.
\end{proof}

At three positions the exterior lengths are $1,7,12,6$.
In degree two there are three regular constituents, two on
each pair diagonal and three on the total diagonal. The
generic rank is three, while the differential-module length
is twelve.

For a constant $d$-dimensional space $V_0$, with $d\geq1$,
write $m^{(d)}_{\pi,q}$ for the composition multiplicity of
$\delta_\pi$ in $\Lambda^q(V_0\otimes H_n)$ and put
$y=(1+t)^d-1$. The same evaluation and inversion give
\begin{align}
 \sum_qm^{(d)}_{\pi,q}t^q
   &=(1+t)^{d s_1(\pi)}
       \prod_{\substack{B\in\pi\\|B|\geq2}}Q_{|B|}(y),
                                      \label{cbc:ct:rank-multiplicity}\\
 \sum_q\operatorname{length}\Lambda^q(V_0\otimes H_n)t^q
   &=\sum_{r=0}^n r!S(n+1,r+1)y^r.\notag
\end{align}
Indeed its exterior series evaluates at $\tau$ to
$(1+t)^{d|\tau|}$. Flatness and generic rank give vanishing
above $dn$. These are also the polynomial Heisenberg
cohomology modules when the constant level form is invertible.
The form is an invertible change of variables in the primitive
differential of Theorem~\ref{cbc:hc:polynomial-cohomology}.

\begin{proposition}[Resolutions and the determinant]
\label{cbc:ct:exterior-resolutions}
A flat differential-module resolution of $\Lambda^qH_n$ has
terms
\[
 E_q^{-a}=\operatorname{Sym}^a(M_n)\otimes_{R_n}
                     \Lambda^{q-a}(M_n),\qquad 0\leq a\leq q,
\]
with differential, for $x_i,y_j\in M_n$ and $b=q-a$,
\begin{align*}
 d(x_1\cdots x_a\otimes y_1\wedge\cdots\wedge y_b)
   =\sum_{j=1}^a x_1\cdots\widehat{x_j}\cdots x_a
       \otimes\partial_wx_j\wedge y_1\wedge\cdots\wedge y_b.
\end{align*}
Its augmentation is the quotient in degree zero. There is also
an ordinary free resolution
\begin{equation}\label{cbc:ct:exterior-telescope}
 0\longrightarrow\bigoplus_m\Lambda^qF_m
 \xrightarrow{\,1-\operatorname{shift}(\Lambda^qT_m)\,}
 \bigoplus_m\Lambda^qF_m\longrightarrow\Lambda^qH_n
 \longrightarrow0,
\end{equation}
using the free stages of Theorem~\ref{cbc:hc:stages}.
In residue coordinates, with
$\Delta_n=\prod_{i<j}(z_i-z_j)$,
\begin{equation}\label{cbc:ct:determinant}
 \Lambda^nH_n\simeq R_n[\Delta_n^{-1}],\qquad
 \Lambda^{dn}(V_0\otimes H_n)
    \simeq R_n[\Delta_n^{-1}]\otimes\det(V_0)^{\otimes n}.
\end{equation}
These last maps are differential-module isomorphisms. Each
module has length $n!$, with multiplicity
$\prod_{B\in\pi}(|B|-1)!$ on $\Delta_\pi$.
\end{proposition}

\begin{proof}
The flat complex $[M_n\xrightarrow{\partial_w}M_n]$ in
degrees $-1,0$ resolves $H_n$. Alternating its $q$-fold
tensor gives $E_q$. The negative-degree variables become
symmetric, and the differential is the displayed sum. Two
replacements cancel by their exterior signs. Since $H_n$ is
flat, the alternating idempotent preserves the tensor-power
quasi-isomorphism, proving exactness.

Exterior powers commute with filtered colimits, so the
cokernel in \eqref{cbc:ct:exterior-telescope} is the required
module. The least occupied input stage of a nonzero finite
sum survives under $1-\operatorname{shift}$, proving
injectivity. Its matrix entries are the corresponding minors
of $T_m$. This is an ordinary-module resolution: a position
derivative can increase the pole stage. The first resolution
itself retains the derivative connections, coefficient base
change and the block decompositions.

For the top exterior power, every transition determinant is a
nonzero scalar times $\Delta_n^2$ by
\eqref{cbc:hc:transition-determinant}. After scalar rescaling,
its colimit is the localization at $\Delta_n$. In generic
residue coordinates the connection is ordinary differentiation.
This proves the first isomorphism in \eqref{cbc:ct:determinant}.
The determinant identity for $V_0\otimes H_n$ proves the
second.

Permuting labels acts on the residue determinant with
its permutation sign, and on the second determinant with that
sign to the power $d$. The top coefficients of
\eqref{cbc:ct:multiplicity} and
\eqref{cbc:ct:rank-multiplicity} give the asserted multiplicities.
Their sum counts permutations by their cycle partitions, and
is $n!$.
\end{proof}

\section{The support left by every consecutive cut}
\label{cbc:ct:support}

The same partition geometry determines the range of the local
vacuum coefficient. For this section $X$ may be any smooth
geometrically connected curve over $k$. Recall
\[
 Z_n=\bigcap_{c=1}^{n-1}\bigcup_{i\leq c<j}\Delta_{ij}
\]
from \eqref{eq:words-support}. We use its reduced closed
structure. Local cohomology depends only on its support. An
equality partition lies in $Z_n$ precisely when each gap is
crossed by one of its blocks.

A \emph{crossing chain} of length $r$ consists of pair blocks
$\{i_a,j_a\}$ and singleton blocks at the other positions,
where
\begin{equation}\label{cbc:ct:crossing-chain}
 1=i_1<i_2<j_1<i_3<j_2<\cdots<i_r<j_{r-1}<j_r=n.
\end{equation}
For $r=1$ this means just the pair $\{1,n\}$.

\begin{proposition}\label{cbc:ct:components}
The irreducible components of $Z_n$ are exactly the diagonals
of crossing chains. There are $\binom{n-2}{2r-2}$ components
of codimension $r$, for $1\leq r\leq\lfloor n/2\rfloor$.
For $n\geq3$ their total number is $2^{n-3}$.
\end{proposition}

\begin{proof}
A component is a diagonal of a refinement-minimal
indecomposable partition. A block crosses exactly the gaps
between its least and greatest labels. Replacing that block by
the pair of these endpoints and making its other labels
singletons preserves the crossed gaps. Minimality therefore
forces pair and singleton blocks.

The intervals $[i_a,j_a-1]$ of its pair blocks cover all gaps,
irredundantly. In order of increasing $i_a$, the $j_a$ also
increase, since otherwise an interval is contained in another.
The next interval must start before the previous pair ends:
$i_{a+1}<j_a$. Equality is impossible because pair blocks have
disjoint endpoints. Irredundancy of the middle interval then
gives $j_a<i_{a+2}$.

The endpoints of the cover are $1,n$.
These are precisely \eqref{cbc:ct:crossing-chain}, and the
same inequalities prove conversely that every chain is an
irredundant cover.

Each pair contributes one independent normal equation. A chain
is determined by its $2r-2$ internal endpoints, whose order
fixes their roles. This proves the codimension and count.
Summing the even binomial coefficients gives $2^{n-3}$ for
$n\geq3$.
\end{proof}

On an affine-line chart one can compute the full supported
object without splitting it into components. Put
$f_c=\prod_{i\leq c<j}(z_i-z_j)$. For any complex $M$ of
coefficients,
\begin{equation}\label{cbc:ct:cut-cech}
 R\Gamma_{Z_n}M\simeq
 M\otimes_{R_n}^{\mathrm L}
       \bigotimes_{c=1}^{n-1}[R_n\longrightarrow R_n[f_c^{-1}]].
\end{equation}

Each bracket has degrees zero and one. The tensor computes
successive supports on $V(f_c)$, hence their intersection.
Its signs use the displayed gap order. It is a flat model and
commutes with ordinary derived coefficient base change.

Let
\[
 \mathcal U_n=R\Gamma_{Z_n}(\omega_{X^n})[n-1]
\]
be the positive coefficient of the cut cobar of the local vacuum.
This is Theorem~\ref{thm:words-support} applied to
$\mathcal J_n=\omega_{X^n}[n]$. For a nonzero bounded complex
$M$, define its width to be the largest nonvanishing cohomological
degree minus the smallest.
Write $i_{1n}:X^{n-1}\hookrightarrow X^n$ for the diagonal
identifying the first and last positions.

\begin{theorem}[Vacuum amplitude]\label{cbc:ct:amplitude}
For $n\geq2$, the lowest cohomology of $\mathcal U_n$ is
\[
 H^{2-n}(\mathcal U_n)=i_{1n,*}\omega_{X^{n-1}}.
\]
For every $1\leq c\leq\lfloor n/2\rfloor$,
$H^{c+1-n}(\mathcal U_n)\ne0$. Thus
\[
 \operatorname{width}(\mathcal U_n)\geq\lfloor n/2\rfloor-1.
\]
For $2\leq n\leq9$ equality holds, and
$H^j(\mathcal U_n)=0$ for $j>\lfloor n/2\rfloor+1-n$.
\end{theorem}

\begin{proof}
Take $V=k$ and $T=0$ in the residue construction. Its maps
are scalar principal parts. The resulting partition complex
resolves $\omega_{X^n}[n]$. To see this, filter $X^n$ by
the closed unions of collision strata of increasing codimension.
The successive localization triangles have, on each stratum
of codimension $c$, only its degree-$c$ normal principal-part
module.

Its extension from the distinct-root locus is exactly
\eqref{cbc:ct:partition}. Locally this open inclusion is a
principal localization and has no higher direct images. The
boundary between successive strata is the oriented binary
principal part.

After the common shift $[n]$, totalizing this
finite filtration gives \eqref{cbc:ct:coalgebra} and resolves
$\omega_{X^n}[n]$. Applying the supported-cobar theorem and
the finite contraction gives $G_n$ for $\mathcal U_n$.

The terms whose partition joins $1$ to $n$ form a subcomplex.
They are the complete scalar Cousin resolution on $\Delta_{1n}$,
shifted by $n-2$, since every further coarsening remains
indecomposable. There is only one indecomposable partition
with $n-1$ blocks, namely $\{1,n\}$ and all other labels
single. Thus the quotient begins in degree $3-n$. The long
exact sequence proves the lowest-degree formula and vanishing
below it.

At the generic point of a crossing-chain component of
codimension $c$, no other component remains. Its local
cohomology is the nonzero top principal-part module in degree
$c$. Chains exist for every stated $c$. After the shift this
proves the nonvanishing and the width bound.

For the upper bound in the finite range, form the poset
$\mathcal P_n$ of intersections of nonempty families of
components, ordered by inclusion of their diagonals. For
$L\in\mathcal P_n$, let $c(L)=\operatorname{codim}L$ and
let $\Delta_L$ be the order complex of the diagonals strictly
containing $L$. The local-cohomology filtration for a linear
arrangement has graded pieces
\begin{equation}\label{cbc:ct:arrangement-filtration}
 H^{c(L)}_L(R_n)\otimes_k
          \widetilde H_{c(L)-r-1}(\Delta_L;k)
          \quad\text{in }H^r_{Z_n}(R_n).
\end{equation}
This is \cite{AMGLZ2003}*{Corollary~1.3}. Tensoring with
the canonical bundle gives the right-module convention.

The following finite calculation determines the needed
vanishing. Enumerate the chains of each strict interval. If
$N_q$ is the number of its $q$-simplices and $\rho_q$ the
rank of the boundary in characteristic two, put
$b_q=N_q-\rho_q-\rho_{q+1}$. Use $N_{-1}=1$, $\rho_{-1}=0$
and the reduced augmentation: $\rho_0$ is one for a nonempty
complex and zero for the empty one. Thus the empty complex
has $b_{-1}=1$.

Let
$B_{n,r}=\sum_L b_{c(L)-r-1}(L)$. The boundary which deletes
one vertex gives
\begin{center}
\begin{tabular}{r|r|rrrr}
\toprule
$n$&$|\mathcal P_n|$&$B_{n,1}$&$B_{n,2}$&$B_{n,3}$&$B_{n,4}$\\
\midrule
2&1&1&0&0&0\\
3&1&1&0&0&0\\
4&3&1&2&0&0\\
5&10&1&10&0&0\\
6&34&1&34&6&0\\
7&126&1&98&78&0\\
8&523&1&258&660&24\\
9&2430&1&642&4764&600\\
\bottomrule
\end{tabular}
\end{center}
All other $B_{n,r}$ are zero. Boundary matrices have integer
entries. Their ranks over characteristic zero are at least
their ranks modulo two. Hence each reduced homology dimension
in characteristic zero is at most the corresponding $b_q$.
The table and \eqref{cbc:ct:arrangement-filtration} give
$H^r_{Z_n}=0$ for $r>\lfloor n/2\rfloor$ in this range.

\'Etale coordinates at the collision strata identify the
diagonal arrangement locally with the linear one. Local
cohomology commutes with this flat pullback, giving the same
vanishing on $X^n$. Shifting by $[n-1]$ and using the
nonvanishing already proved completes the argument.
\end{proof}

The table gives dimensions in characteristic two, used only
as bounds in characteristic zero. It does not split the
filtration. The lower bound holds in every arity, so no
choice of shifts makes all vacuum coefficients have a
uniformly bounded width. The corresponding upper bound for
$n\geq10$ remains open.
